%======================================================================
\subsection{BW-SIM-06 plan: airflow with OpenFOAM (31 Oct 2026)}\label{sec:sim06-plan}
%======================================================================

You asked to start \eref{SIM-06}, which had been left for later. This section sets the scope, the theory and
the checks against known answers, and ends with your questions. No code is written until you answer them.

\paragraph{What it is for.} For SP-01 and the parts around it, airflow answers four questions that the
structural and thermal studies cannot:
\begin{itemize}[nosep]
  \item \textbf{Prop wash.} How much the propeller wake pushes down on the arms and frame (\emph{download},
        which costs thrust), and the load it puts on them.
  \item \textbf{Frame drag} in forward flight at a given speed and pitch.
  \item \textbf{Cooling.} How much air reaches the ESC, motors and battery, and what heat-transfer coefficient
        $h$ it gives on each surface.
  \item \textbf{Feeding the thermal study.} That $h$ replaces the estimated convection presets (still air 10,
        light airflow 25, prop wash 75\,W/m$^2$K) of \S\ref{sec:sim-thermal-plan}.
\end{itemize}
Out of scope: designing propeller blades (modelling the spinning blades themselves), noise, and unsteady
vortex shedding.

\paragraph{Why not the CfdOF add-on.} The roadmap named CfdOF, the FreeCAD add-on that drives OpenFOAM. It
lists macOS as ``not widely tested'' and points Mac users to a Docker container. It also leans on cfMesh and
HiSA, which the native Apple-silicon OpenFOAM build does not ship. Vector Labs already drives CalculiX
directly from its own study frame (\S\ref{sec:sim01-built}), so CFD follows the same pattern: the study writes
an OpenFOAM case folder, runs the solver as a separate program, and reads the results back. CfdOF is LGPL and
its dictionary templates can be borrowed where useful.

\paragraph{OpenFOAM on this Mac.} The \emph{openfoam-app} project (G.~Gerlero) packages the official ESI
OpenFOAM releases as a native macOS app: OpenFOAM-v2606, Apple silicon, macOS 14 or later, a 268\,MB
download. It is installable from Homebrew or as a zip, and is started with
\file{OpenFOAM-v2606.app/Contents/Resources/etc/openfoam}. It includes everything this plan uses:
\code{blockMesh}, \code{snappyHexMesh}, \code{simpleFoam}, the actuator-disk source, \code{foamToVTK}, and MPI
for parallel runs. OpenFOAM is GPL-3. Like CalculiX, it runs as a separate program, so Vector Labs' own
licence is unaffected. Bundling it means shipping its licence and a pointer to its source (question O1).

%----------------------------------------------------------------------
\subsubsection*{Theory}

\paragraph{The flow equations.} A drone flies far below Mach~0.3, so the air is treated as incompressible.
The steady Reynolds-averaged Navier--Stokes (RANS) equations for the mean velocity $\mathbf u$ and pressure $p$
are
\[
  \nabla\cdot\mathbf u=0,\qquad
  (\mathbf u\cdot\nabla)\mathbf u=-\frac{1}{\rho}\nabla p
   +\nabla\cdot\bigl[(\nu+\nu_t)(\nabla\mathbf u+\nabla\mathbf u^{\mathsf T})\bigr]+\frac{\mathbf f}{\rho},
\]
with $\nu$ the air's viscosity and $\nu_t$ the turbulent viscosity. $\nu_t$ comes from the $k$--$\omega$ SST
model, the usual choice for flow around bodies that separates. $\mathbf f$ is a body force per unit volume,
used for the propellers below. \code{simpleFoam} solves these equations with the SIMPLE pressure--velocity
iteration until the residuals and the force on the body stop changing. Air properties come from the study's
temperature and altitude (standard atmosphere): $\rho=1.204$\,kg/m$^3$ and $\nu=1.51\times10^{-5}$\,m$^2$/s
at 20\,$^\circ$C at sea level.

\paragraph{Propellers as actuator disks.} The spinning blades are not modelled. Each propeller becomes a thin
disk of area $A=\pi D^2/4$ and thickness $t$ that adds momentum to the air. By momentum theory, a disk with
thrust $T$ in hover accelerates the air through it to the induced velocity, and the far wake to twice that:
\[
  v_i=\sqrt{\frac{T}{2\rho A}},\qquad v_{\text{wake}}=2v_i,\qquad
  \mathbf f=\frac{T}{A\,t}\,\hat{\mathbf n}\quad\text{inside the disk.}
\]
In forward flight at speed $V$, $T=2\rho A\,v_i\,\lvert\mathbf V+v_i\hat{\mathbf n}\rvert$ is solved for $v_i$.
The thrust per motor comes from the \eref{ROB-04} thrust table at the hover throttle, which the
Mass \& Balance panel already computes (\S\ref{sec:rob0304-plan}), so the wake matches the drone's real
weight. Swirl, a tangential force from the motor torque $Q$, is an option. OpenFOAM provides the disk as the
\code{actuationDiskSource} option on a cell zone made from the motor's \code{ThrustAxis}.

\paragraph{Heat transfer from the flow.} For cooling, the first step is one-way and cheap. The flow is solved
with the heated surfaces held at a reference temperature $T_w$. The solver reports the wall heat flux $q''$
on each face, and the study turns it into a surface coefficient
\[
  h=\frac{q''}{T_w-T_\infty}.
\]
These per-face $h$ values are written out as a \emph{convection map}, which the thermal study reads in place
of a preset. The temperature of solid and air solved together (conjugate heat transfer, OpenFOAM's
\code{chtMultiRegionFoam}) is more exact but much slower, and is left for later (question O2).

\paragraph{Mesh.} \code{blockMesh} fills a box of hexahedral cells sized from the model's bounding box; the
box and refinements are sketched in Fig.~\ref{fig:sim06-domain}. \code{snappyHexMesh} then cuts the model out
of the box. The model is exported as STL, one file per part, so forces can be reported per part. Cells are
refined near the walls and inside each propeller disk and its wake, and thin cell layers are added along the
walls. On the walls, wall functions are used: the first cell sits at a dimensionless wall distance
$y^+\approx30$--$300$, which needs far fewer cells than resolving the boundary layer. The pre-check estimates
$y^+$ from the speed and the cell size. Presets: coarse about 0.3\,M cells, medium 1\,M, fine 3\,M, run in
parallel on the performance cores. Run times will be measured on the benchmarks, not promised here.

\begin{figure}[ht]
\centering
\begin{tikzpicture}[scale=0.85,font=\small]
  % domain
  \draw[thick] (-5,-4.6) rectangle (5,2.2);
  \node[anchor=north west] at (-5,2.2) {domain (from the bounding box)};
  % refinement box
  \draw[dashed,bwdark] (-3.3,-3.8) rectangle (3.3,1.0);
  \node[bwdark,anchor=south east] at (3.3,1.0) {refinement};
  % frame
  \fill[gray!60] (-2.4,-0.12) rectangle (2.4,0.12);
  \fill[gray!80] (-0.7,-0.35) rectangle (0.7,0.35);
  % disks
  \foreach \x in {-2.1,2.1}{
    \fill[red!60] (\x-0.85,0.32) rectangle (\x+0.85,0.42);
    % stream tube contracting
    \draw[blue!60] (\x-0.85,0.37) .. controls (\x-0.75,-1.2) .. (\x-0.6,-3.6);
    \draw[blue!60] (\x+0.85,0.37) .. controls (\x+0.75,-1.2) .. (\x+0.6,-3.6);
    \draw[->,blue!70,thick] (\x,0.1) -- (\x,-1.0) node[right] {$v_i$};
    \draw[->,blue!70,thick] (\x,-2.4) -- (\x,-3.6) node[right] {$2v_i$};
  }
  \node[red!70!black] at (-2.1,0.75) {disk $T$};
  \node at (0,0.75) {frame};
  \draw[->] (-4.6,1.4) -- (-3.8,1.4) node[right] {$\mathbf V$ (forward flight)};
  \node[anchor=south] at (0,-4.6) {outlet};
\end{tikzpicture}
\caption{Prop-wash study in hover. Each propeller is a red actuator disk with thrust $T$ from the thrust
table. The wake contracts and speeds up from $v_i$ at the disk to $2v_i$ downstream, and strikes the arms. The
refinement box keeps the cells fine around the frame and the wake. In forward flight the free stream $\mathbf
V$ enters at the side.}
\label{fig:sim06-domain}
\end{figure}

\paragraph{Convergence.} A run is converged when the residuals fall below $10^{-4}$, the force on each part
changes less than 0.5\,\% over the last 200 iterations, and the mass flow in matches the mass flow out to
within 0.5\,\%. The report states all three. Mesh independence comes from running the coarse and the medium
mesh: when the forces differ by more than 5\,\%, the report says so and suggests the fine mesh.

%----------------------------------------------------------------------
\subsubsection*{The workflow}

New Study gains an \emph{Airflow} group with three types.
\begin{description}[nosep,font=\normalfont\bfseries]
  \item[External flow.] Frame drag at a speed and an angle of attack. Outputs: drag and lift per part, the
        coefficients $C_d$ and $C_l$ on a reference area you choose, surface pressure, and the flow on a plane.
  \item[Prop wash.] Hover or forward flight, with a disk on each motor that has a \code{ThrustAxis} and a
        thrust from the \eref{ROB-04} table. Outputs: the download on each part as a percentage of the total
        thrust, plus the wall pressure on the arms.
  \item[Cooling.] Openings, inlets or fans (a fixed flow rate or a prop-wash disk) and the heated surfaces.
        Outputs: the air speed at each part and the convection map for the thermal study.
\end{description}
\begin{itemize}[nosep]
  \item \textbf{Setup:} air (temperature and altitude), the free stream, the domain size (automatic, which
        you can change) and the mesh preset.
  \item \textbf{Pre-check:} watertight parts, an estimate of the cell count and of $y^+$, and a disk on
        every marked motor.
  \item \textbf{Run:} in the background, with a live residual and force plot and a Stop button.
  \item \textbf{Results:} read through \code{foamToVTK} into FreeCAD's FEM result pipeline: cut planes,
        streamlines, surface pressure and the forces table. When ParaView is installed, an \emph{Open in
        ParaView} button opens the case there.
  \item \textbf{Report:} the \LaTeX{} report gains an airflow section with the setup, convergence and force
        tables.
\end{itemize}

%----------------------------------------------------------------------
\subsubsection*{Benchmarks (as in \eref{SIM-05})}
\begin{center}
\small
\begin{tabularx}{\linewidth}{@{}lX@{}}
\toprule
Pipe, laminar & Air at 20\,$^\circ$C, $D=10$\,mm, mean speed 0.5\,m/s ($Re=331$). Fully developed pressure
  gradient $\mathrm dp/\mathrm dx=32\mu U/D^2=2.91$\,Pa/m, measured past the entrance length
  ($0.05\,Re\,D\approx0.17$\,m); accepts 2\,\%.\\
Sphere drag & $Re=100$. The correlation of Schiller and Naumann, $C_d=\frac{24}{Re}(1+0.15Re^{0.687})=1.09$;
  accepts 5\,\%.\\
Actuator disk & Hover, $T=1$\,N, $D=76$\,mm: $v_i=\sqrt{T/2\rho A}=9.57$\,m/s at the disk (mean across it)
  and 19.1\,m/s in the far wake; accepts 5\,\% at the disk and 10\,\% in the wake (momentum theory ignores
  mixing).\\
Flat plate, heated & Air at 5\,m/s along a 50\,mm plate, laminar: $Re_L=1.66\times10^4$,
  $\bar h=\frac{k}{L}\,0.664\,Re_L^{1/2}Pr^{1/3}=39$\,W/m$^2$K; accepts 10\,\%.\\
\bottomrule
\end{tabularx}
\end{center}

\paragraph{Build order.} Each phase is one Sonnet agent, tested against its benchmark before the next starts.
\begin{enumerate}[nosep]
  \item The case writer and runner, external flow and the results view (pipe and sphere benchmarks).
  \item Actuator disks and the prop-wash study, linked to the thrust tables (disk benchmark).
  \item Cooling and the convection map used by the thermal study (flat-plate benchmark).
  \item Wall pressure as a load on a static study (question O7).
\end{enumerate}

%----------------------------------------------------------------------
\paragraph{Questions for you.}
\begin{enumerate}[label=O\arabic*.]
  \item \textbf{Getting OpenFOAM.} (a) Bundle it inside Vector Labs, like CalculiX and Tectonic: one download,
        always the tested version, but the app grows by about 270\,MB zipped and more unpacked. (b) Use an
        installed \emph{OpenFOAM.app}: the study finds it, and if none is installed it offers an
        \emph{Install OpenFOAM\ldots} button (Homebrew or the download). \emph{Recommendation:} (b), with the
        tested version pinned and a warning on any other version, so the app stays lean.
  \item \textbf{First study types.} Prop wash and external flow first, then cooling with the convection map;
        the full solid-and-air temperature solution (conjugate heat transfer) later.
        \emph{Recommendation:} yes.
  \item \textbf{Propellers.} Actuator disks with thrust from the \eref{ROB-04} thrust table at hover (or a
        typed thrust), and swirl from the torque as an option, off by default. Modelling the spinning blades
        themselves is out of scope. \emph{Recommendation:} yes.
  \item \textbf{Turbulence.} $k$--$\omega$ SST by default, laminar as an option for slow, small flows (the pre-check
        suggests it below $Re\approx2000$). \emph{Recommendation:} yes.
  \item \textbf{Mesh.} \code{snappyHexMesh} with the coarse, medium and fine presets and an automatic domain,
        all editable. \emph{Recommendation:} yes.
  \item \textbf{Results.} Inside Vector Labs (planes, streamlines, surface pressure, forces table), plus
        \emph{Open in ParaView} when it is installed. \emph{Recommendation:} yes.
  \item \textbf{Linking to the other studies.} (a) The convection map feeds the thermal study. (b) The wall
        pressure becomes a load on a static study (for example, the arms under prop wash).
        \emph{Recommendation:} both, (a) first.
  \item \textbf{Where cases live.} The OpenFOAM case folders go in \file{\textasciitilde/Library/Caches/Vector
        Labs/cfd/}, one per study, kept until you clear them. The document stores only the setup and the
        result summary; a re-run rebuilds the case. Runs use all performance cores.
        \emph{Recommendation:} yes.
  \item \textbf{Acceptance.} The four benchmarks above, each run in CI at the coarse preset.
        \emph{Recommendation:} yes.
\end{enumerate}

\paragraph{Your answers (31 Oct).} O1--O9 as recommended. OpenFOAM is \emph{not} bundled: Vector Labs finds
an installed \emph{OpenFOAM-v2606.app} (the tested version; any other version triggers a warning) and offers
\emph{Install OpenFOAM\ldots} when none is installed. Prop wash and external flow come first, then cooling with
the convection map. Propellers are actuator disks fed by the thrust tables, with swirl off by default.
Turbulence is $k$--$\omega$ SST, with laminar as an option. The mesh uses \code{snappyHexMesh} presets with an
automatic domain. Results show inside Vector Labs, with \emph{Open in ParaView} when it is installed. The links
to the other studies are the convection map (built first) and wall pressure as a static load. Cases live in
\file{\textasciitilde/Library/Caches/Vector Labs/cfd/}. The four benchmarks are the acceptance, and they will
be joined by the open validation data now being gathered.
